\chapter*{Abstract}
\addcontentsline{toc}{chapter}{Abstract}

Swype Hero is a browser-based arcade game in which descending balloons are destroyed by drawing the geometric shapes they carry. The game is designed to be enjoyable on its own terms, and it is also designed as a curriculum: a sequence of gestures of increasing geometric and motor difficulty that eventually converges on the trajectories traced by words across a touchscreen keyboard. This document is both a game design and a technical specification. It defines the playfield, the balloon dynamics, the gesture lifecycle, the representation and geometry of gestures, a family of recognition algorithms with explicit confidence and rejection behavior, a learning and difficulty model, the construction of word trajectories from keyboard geometry, the software architecture of a browser implementation, a roadmap of later input devices, and an experimental research program for studying gesture recognition and motor learning with consent-based datasets.

Three commitments organize the whole. Recognition is independent of gameplay, so that the same engine can serve a typing tutor, a keyboard, or another application. Raw input is kept separate from interpreted gestures, so that recognition errors can be investigated and algorithms compared without rewriting recorded observations. Learning by the player and adaptation by the software are treated as distinct processes that must remain experimentally distinguishable. The specification states normative requirements with the conventions described below, marks the evidential status of its historical and behavioral claims, and defers every empirical question it cannot settle to a named hypothesis and an experimental protocol.

\chapter*{Motivation}
\addcontentsline{toc}{chapter}{Motivation}

Swype Hero is motivated by the observation that gesture-based text entry remains an underexplored form of human-computer interaction, despite the remarkable proficiency that users can develop through prolonged practice. The original Swype keyboard demonstrated that continuous movements across a keyboard could become a highly efficient means of producing language. Experienced users learned not merely to trace letters, but to develop distinctive motor strategies involving curves, sharp turns, abbreviated trajectories, and deliberate exaggerations that made otherwise ambiguous words recognizable. Over time, the relationship between the user and the recognition system could become increasingly specialized.

The discontinuation of Swype also illustrates a broader problem in software development: improvements in contemporary products do not necessarily preserve the capabilities of their predecessors. Gesture shortcuts, personalized dictionaries, and recognition behaviors developed through years of use can disappear when an application is abandoned or replaced. The resulting loss is not merely functional. It can invalidate a substantial investment in acquired motor proficiency.

Swype Hero approaches this problem through play. Inspired in part by gesture-drawing games such as Magic Touch, it begins with a simple premise: players destroy balloons by reproducing increasingly complex geometric shapes. What initially appears to be an arcade mechanic gradually becomes a curriculum in continuous gesture production, progressing from lines, angles, and loops to intricate trajectories resembling the shapes of written words. Eventually, the game can introduce actual keyboard-derived word gestures, transforming accumulated motor skills into practical text-entry proficiency.

The project also provides an experimental environment for investigating gesture recognition, motor learning, ambiguity, and personalization. Rather than assuming that users should reproduce a single ideal trajectory, it asks how recognition systems might accommodate the distinctive movements that individuals develop through practice. Its longer-term architecture anticipates alternative input methods, including mice, keyboards, joysticks, Etch A Sketch-style controls, and eye tracking, allowing a common gesture vocabulary to be explored across different physical interfaces.

The immediate objective is modest: to create an enjoyable, responsive browser game with progressively challenging gestures. The larger ambition is to investigate how a playful interaction can become a general method of acquiring computational fluency, and how software might preserve and extend that fluency rather than repeatedly requiring users to begin again.

\chapter*{Conventions of This Specification}
\addcontentsline{toc}{chapter}{Conventions of This Specification}

\paragraph{Normative language.} The words \must, \mustnot, \should, \shouldnot{} and \may, set in small capitals, are used in the sense of RFC~2119 \cite{rfc2119}. Statements containing them appear inside numbered requirements. Text outside a numbered requirement is explanatory unless it says otherwise. A requirement binds an implementation only within the milestone that introduces it (Part~\ref{part:roadmap}); later parts describe designs for later milestones and are normative only when those milestones are adopted.

\paragraph{Evidential status.} The project draws on published research, on the documented behavior of software, on the reported experience of long-term users, and on its own design decisions. These are different kinds of claim and the text marks them. A superscript \evD{} marks a claim documented in the published literature or in primary technical documentation. A superscript \evO{} marks a behavioral observation reported from practice that has not been independently verified. A superscript \evH{} marks a hypothesis that this project proposes to test, and each such hypothesis is stated as a numbered hypothesis with a protocol in Part~\ref{part:research}. A superscript \evS{} marks a design decision of this specification, which is a commitment and not a claim about the world. A superscript \evC{} marks a constant that awaits calibration: the requirement that uses it is executable and testable once a value is supplied, but the value, its unit and the distance to which it applies are fixed by a stated procedure (Appendix~\ref{app:notation}) and not by this text. A requirement without that mark is testable as written. Unmarked statements are definitions, arithmetic consequences of definitions, or connective prose.

\paragraph{Coordinates and units.} The playfield uses logical units defined in Chapter~\ref{ch:playfield}, with the origin at the top-left corner and the $y$ axis pointing down, as in the canvas coordinate systems of browsers. Angles are in radians unless stated otherwise. Distances on the keyboard are measured in key widths. Times are in milliseconds and are taken from a monotonic clock.

\paragraph{Computed values.} Every numerical value reported as computed in this document, including the word-geometry figure and tables of Part~\ref{part:words}, is produced by the script \texttt{tools/wordgeom.py} distributed with the source (and the edge-case fixtures of Appendix~\ref{app:fixtures} by \texttt{tools/fixtures\_check.py}), from the definitions given in the text. Numerical constants chosen as defaults, such as thresholds and budgets, are starting points for calibration and are labeled as design decisions.

\paragraph{Notation.} A gesture is written $G$ and a library template $T$. The number of points after resampling is $N$. The recognizer's base distance is $d$. Three derived quantities are kept apart: the quotient pseudometric $d_H$ under a common isometry group $H$, the directional matching cost $c_g$ of a stroke against a candidate gesture $g$ under its allowed set $\mathcal S_g$, and the symmetric confusability $\kappa$ between templates of distinct gestures (Chapter~\ref{ch:invariance}). A fuller table of symbols is given in Appendix~\ref{app:notation}.

\paragraph{Structure.} Parts~I and~II specify the concept and the playable game. Part~III develops the geometry on which everything else rests, and Part~IV the algorithms that act on it. Part~V treats learning, Part~VI the passage from shapes to words, and Part~VII presentation. Part~VIII is the software architecture. Part~IX anticipates other input devices, Part~X is the research program, Part~XI describes applications beyond the game, and Part~XII sets out the roadmap and the scope of the first implementation. Appendices collect schemas, reference algorithms, protocols and the historical record.
